If $\varkappa_1, \ldots, \varkappa_n$ are $n$ positive quantities such that

$$\varkappa_1 \cdot \varkappa_2 \cdots \varkappa_n = |\sqrt{d}|,$$

$$\varkappa_{p+r_2} = \varkappa_p \quad \text{for } p = r_1 + 1, \ldots, r_1 + r_2,$$

then by Theorem 95 there is a nonzero integer $\alpha$ in $K$ (whose norm thus has at least absolute value 1) such that

$$|\alpha^{(i)}| \leq \varkappa_i \quad \text{for } i = 1, 2, \ldots, n, 1 \leq |N(\alpha)| \leq \sqrt{d}.$$

From this it follows that

$$|\alpha^{(i)}| \geq \frac{1}{|\alpha^{(1)}| \cdot |\alpha^{(2)}| \cdots |\alpha^{(i-1)}| \cdot |\alpha^{(i

and for the corresponding $\alpha_h$ we have

$$\left| L(\alpha_h) - 2Ah \right| < A$$
$$A(2h - 1) < L(\alpha_h) < A(2h + 1).$$

Consequently

$$L(\alpha_1) < L(\alpha_2) < L(\alpha_3) \ldots$$

(62)

while at the same time

$$\left| N(\alpha_h) \right| \leq \left| \sqrt{d} \right|.$$

These infinitely many principal ideals $(\alpha_h)$ whose norms are not greater than $\sqrt{d}$, cannot all be distinct; hence there must exist at least two distinct indices $h$ and $m$ such that

$$(\alpha_h) = (\alpha_m),$$ hence $\alpha_m = \varepsilon\alpha_h,$

where $\varepsilon$ is a unit in $k$. For this $\varepsilon$, by (62),

$$L(\alpha_h) \neq L(\alpha_m) = L(\varepsilon\alpha_h)$$
$$L(\varepsilon) = L(\varepsilon\alpha_h) - L(\alpha_h) \neq 0,$$

and Lemma (b) is proved. From this we obtain

**Lemma (c).** If $r \geq 0$, then the number $k$ of independent units of the field is exactly $r$.

By (b) there is a unit $\varepsilon_1$ such that

$$\log \left| \varepsilon_1^{(1)} \right| \neq 0.$$

Then if $r > 1$, there is likewise a unit $\varepsilon_2$ such that

$$\begin{vmatrix} \log \left| \varepsilon_1^{(1)} \right| & \log \left| \varepsilon_2^{(1)} \right| \\ \log \left| \varepsilon_1^{(2)} \right| & \log \left| \varepsilon_2^{(2)} \right| \end{vmatrix} \neq 0$$

and so on. Thus we conclude from (b) the existence of $r$ units $\varepsilon_1, \ldots, \varepsilon_r$ for which the determinant

$$

order is exactly $r$, and together with Lemma (d), the validity of Dirichlet’s unit theorem, Theorem 100, is established.

By Theorem 38, §11, we see at once that between two systems of units in $K: \eta_1, \ldots, \eta_r$ and $\varepsilon_1, \ldots, \varepsilon_r$, equations of the form

$$\eta_m = \zeta_m \varepsilon_1^{a_{m1}} \cdot \varepsilon_2^{a_{m2}} \cdots \varepsilon_r^{a_{mr}} \quad (m = 1, 2, \ldots, r)$$

hold, where the $\zeta_m$ are roots of unity, while the $a_{mk}$ are rational integers with determinant $\pm 1$. Thus for each system of fundamental units of $K$ the absolute value of the determinant

$$\begin{vmatrix}
\log |\eta_1^{(1)}| & \cdots & \log |\eta_r^{(1)}| \\
\vdots & & \vdots \\
\log |\eta_1^{(r)}| & \cdots & \log |\eta_r^{(r)}|
\end{vmatrix} = \pm R$$

has the same nonzero value; thus this value is a constant of the field.

The absolute value $R$ of the determinant

$$\begin{vmatrix}
e_1 \log |\eta_1^{(1)}| & \cdots & e_1 \log |\eta_r^{(1)}| \\
\vdots & & \vdots \\
e_r \log |\eta_1^{(r)}| & \cdots & e_r \log |\eta_r^{(r)}|
\end{vmatrix} = \pm R$$

is called the regulator $R$ of the field $K$.
